\documentclass[12pt, a4paper]{article}

\usepackage[utf8]{utf8}
\usepackage[margin=1in]{geometry}
\usepackage{graphicx}
\usepackage{hyperref}
\usepackage{listings}
\usepackage{xcolor}
\usepackage{booktabs}
\usepackage{amsmath}
\usepackage{enumitem}
\usepackage{titlesec}
\usepackage{caption}

\definecolor{codegreen}{rgb}{0,0.6,0}
\definecolor{codegray}{rgb}{0.5,0.5,0.5}
\definecolor{codepurple}{rgb}{0.58,0,0.82}
\definecolor{backcolour}{rgb}{0.96,0.96,0.96}
\definecolor{primaryblue}{rgb}{0.1,0.3,0.6}

\hypersetup{
    colorlinks=true,
    linkcolor=primaryblue,
    filecolor=magenta,      
    urlcolor=cyan,
    pdftitle={Smart Hospital Management System - Project Report},
}

\lstdefinestyle{customstyle}{
    backgroundcolor=\color{backcolour},   
    commentstyle=\color{codegreen},
    keywordstyle=\color{primaryblue}\bfseries,
    numberstyle=\tiny\color{codegray},
    stringstyle=\color{codepurple},
    basicstyle=\ttfamily\footnotesize,
    breakatwhitespace=false,         
    breaklines=true,                 
    captionpos=b,                    
    keepspaces=true,                 
    numbers=left,                    
    numbersep=5pt,                  
    showspaces=false,                
    showstringspaces=false,
    showtabs=false,                  
    tabsize=2
}
\lstset{style=customstyle}

\titleformat{\section}{\large\bfseries\color{primaryblue}}{\thesection}{1em}{}
\titleformat{\subsection}{\normalfont\bfseries\color{primaryblue}}{\thesubsection}{1em}{}

\title{
    \textbf{\huge Smart Hospital Management System}\\
    \Large Database System \& Dynamic Web Application Project Report\\
    \vspace{0.5cm}
    \small Prepared for Academic Presentation \& Overleaf Documentation
}
\author{
    \textbf{Abu Fahad Biddut}\\
    Department of Computer Science \& Engineering
}
\date{\today}

\begin{document}

\maketitle

\begin{abstract}
This document provides a comprehensive project report for the \textbf{Smart Hospital Management System (SHMS)}, a dynamic web application built using PHP, MySQL, HTML5, CSS3, and JavaScript. The system is designed to streamline administrative workflows, patient registrations, doctor schedules, department allocations, room bookings, in-patient admissions, inventory control, and automated billing operations. This report outlines the database design, Entity-Relationship (ER) schema, Data Definition Language (DDL) and Data Manipulation Language (DML) scripts, core backend implementation, dynamic CRUD operations for Room, Admission, and Billing modules, and integration with foundational laboratory exercises.
\end{abstract}

\tableofcontents
\newpage

\section{Introduction \& System Overview}
Healthcare facilities require robust, data-driven management tools to maintain seamless operations between medical practitioners, administrative staff, and patients. The \textbf{Smart Hospital Management System} addresses these needs by centralizing operational workflows into a relational database model.

\subsection{Key Objectives}
\begin{itemize}
    \item \textbf{Centralized User Authentication}: Secure access for Administrators and registered Users.
    \item \textbf{Department \& Medical Specialist Management}: Categorization of departments (e.g., Cardiology, Neurology, Pediatrics) and physician details.
    \item \textbf{Patient Record Tracking}: Detailed logging of demographics, blood groups, contact information, and medical histories.
    \item \textbf{Appointment Scheduling}: Real-time booking between active patients and designated specialists.
    \item \textbf{Room \& Ward Logistics}: Room inventory management across General, Private, ICU, and VIP categories with status tracking (Available/Occupied).
    \item \textbf{In-Patient Admissions}: Tracking patient check-ins, assigned rooms, admission dates, and automated room status updates.
    \item \textbf{Integrated Automated Billing}: Automatic invoice generation calculating room tariffs based on days stayed combined with pharmacy expenses.
\end{itemize}

\section{Database Architecture \& Schema Design}
The relational schema for \texttt{hospital\_db} consists of nine interconnected tables enforced by foreign key constraints and cascading operations to maintain transactional integrity.

\subsection{Relational Schema Overview}
\begin{enumerate}
    \item \textbf{\texttt{users}}: Stores user credentials, email addresses, and system roles (\texttt{Admin}, \texttt{User}).
    \item \textbf{\texttt{departments}}: Maintains hospital departments and clinical domain descriptions.
    \item \textbf{\texttt{doctors}}: Contains medical staff data linked to a specific department via \texttt{department\_id}.
    \item \textbf{\texttt{patients}}: Stores patient demographic data (name, gender, age, blood group, address).
    \item \textbf{\texttt{appointments}}: Maps scheduled appointments between patients and doctors with date/time stamps.
    \item \textbf{\texttt{rooms}}: Tracks floor levels, room types (\texttt{General}, \texttt{Private}, \texttt{ICU}, \texttt{VIP}), availability statuses, and daily charges.
    \item \textbf{\texttt{admissions}}: Records patient check-in/check-out dates associated with assigned rooms.
    \item \textbf{\texttt{medicines}}: Pharmacy inventory tracking unit prices and stock quantities.
    \item \textbf{\texttt{bills}}: Calculates final payment obligations based on room stays and medication costs.
\end{enumerate}

\subsection{Entity Relationship Summary}
\begin{table}[h!]
\centering
\caption{Database Table Relationships \& Foreign Key Mapping}
\begin{tabular}{llll}
\toprule
\textbf{Child Table} & \textbf{Foreign Key} & \textbf{Parent Table} & \textbf{On Delete / Update} \\
\midrule
\texttt{doctors} & \texttt{department\_id} & \texttt{departments(id)} & CASCADE / CASCADE \\
\texttt{appointments} & \texttt{patient\_id} & \texttt{patients(id)} & CASCADE / CASCADE \\
\texttt{appointments} & \texttt{doctor\_id} & \texttt{doctors(id)} & CASCADE / CASCADE \\
\texttt{admissions} & \texttt{patient\_id} & \texttt{patients(id)} & CASCADE / CASCADE \\
\texttt{admissions} & \texttt{room\_id} & \texttt{rooms(id)} & CASCADE / CASCADE \\
\texttt{bills} & \texttt{patient\_id} & \texttt{patients(id)} & CASCADE / CASCADE \\
\texttt{bills} & \texttt{admission\_id} & \texttt{admissions(id)} & CASCADE / CASCADE \\
\bottomrule
\end{tabular}
\end{table}

\section{Database SQL DDL Implementation}
The following SQL script demonstrates the creation of the core tables, foreign keys, and default seed data.

\begin{lstlisting}[language=SQL, caption={Hospital Database SQL Schema (hospital\_db.sql)}]
CREATE TABLE `rooms` (
  `id` INT AUTO_INCREMENT PRIMARY KEY,
  `room_number` VARCHAR(20) NOT NULL UNIQUE,
  `room_type` ENUM('General', 'Private', 'ICU', 'VIP') NOT NULL,
  `floor` INT NOT NULL,
  `status` ENUM('Available', 'Occupied', 'Maintenance') NOT NULL DEFAULT 'Available',
  `charge_per_day` DECIMAL(10,2) NOT NULL DEFAULT 800.00
) ENGINE=InnoDB DEFAULT CHARSET=utf8mb4;

CREATE TABLE `admissions` (
  `id` INT AUTO_INCREMENT PRIMARY KEY,
  `patient_id` INT NOT NULL,
  `room_id` INT NOT NULL,
  `admission_date` DATE NOT NULL,
  `discharge_date` DATE DEFAULT NULL,
  `status` ENUM('Admitted', 'Discharged') NOT NULL DEFAULT 'Admitted',
  CONSTRAINT `fk_admissions_patient` FOREIGN KEY (`patient_id`) REFERENCES `patients` (`id`) ON DELETE CASCADE,
  CONSTRAINT `fk_admissions_room` FOREIGN KEY (`room_id`) REFERENCES `rooms` (`id`) ON DELETE CASCADE
) ENGINE=InnoDB DEFAULT CHARSET=utf8mb4;

CREATE TABLE `bills` (
  `id` INT AUTO_INCREMENT PRIMARY KEY,
  `patient_id` INT NOT NULL,
  `admission_id` INT NOT NULL,
  `medicine_cost` DECIMAL(10,2) NOT NULL DEFAULT 0.00,
  `room_charge` DECIMAL(10,2) NOT NULL DEFAULT 0.00,
  `total_amount` DECIMAL(10,2) NOT NULL DEFAULT 0.00,
  `payment_status` ENUM('Pending', 'Paid') NOT NULL DEFAULT 'Pending',
  CONSTRAINT `fk_bills_patient` FOREIGN KEY (`patient_id`) REFERENCES `patients` (`id`) ON DELETE CASCADE,
  CONSTRAINT `fk_bills_admission` FOREIGN KEY (`admission_id`) REFERENCES `admissions` (`id`) ON DELETE CASCADE
) ENGINE=InnoDB DEFAULT CHARSET=utf8mb4;
\end{lstlisting}

\section{Room, Admission \& Billing Implementation}
The application features specialized modules for room management, patient admissions, and financial invoicing.

\subsection{Patient Discharge \& Automated Billing Logic}
Below is the PHP handler (\texttt{admission/discharge.php}) that calculates total stay length, updates room availability, and generates a bill invoice:

\begin{lstlisting}[language=PHP, caption={Patient Discharge \& Billing Handler (admission/discharge.php)}]
<?php
include "../connection.php";

$admission_id = intval($_GET['id']);
$sql_adm = "SELECT admissions.*, rooms.charge_per_day, rooms.id AS room_id 
            FROM admissions 
            JOIN rooms ON admissions.room_id = rooms.id 
            WHERE admissions.id = $admission_id";
$res_adm = mysqli_query($conn, $sql_adm);
$adm = mysqli_fetch_assoc($res_adm);

$patient_id = $adm['patient_id'];
$room_id = $adm['room_id'];
$days_stayed = max(1, round((strtotime(date('Y-m-d')) - strtotime($adm['admission_date'])) / 86400));
$room_total_charge = $days_stayed * $adm['charge_per_day'];
$medicine_cost = 250.00;
$total_amount = $room_total_charge + $medicine_cost;

mysqli_query($conn, "UPDATE admissions SET discharge_date = CURRENT_DATE(), status = 'Discharged' WHERE id = $admission_id");
mysqli_query($conn, "UPDATE rooms SET status = 'Available' WHERE id = $room_id");

$sql_bill = "INSERT INTO bills (patient_id, admission_id, medicine_cost, room_charge, total_amount, payment_status) 
             VALUES ($patient_id, $admission_id, $medicine_cost, $room_total_charge, $total_amount, 'Pending')";
mysqli_query($conn, $sql_bill);
?>
\end{lstlisting}

\section{Verification \& Conclusion}
The Smart Hospital Management System has been tested across Apache/MySQL environments (XAMPP). All relational integrity constraints, foreign key cascades, dynamic page routing, room allocations, discharge procedures, and SQL billing operations function accurately.

\end{document}
